% RED-1 Research Memo: Intrinsic Homological Characterization of Chiral Koszulness
% Date: 2026-03-17
% Author: RED-1 (pure homological algebraist)

\documentclass[11pt]{article}
\usepackage{amsmath,amssymb,amsthm,enumitem}
\usepackage[margin=1in]{geometry}

\newtheorem{conjecture}{Conjecture}
\newtheorem{proposition}[conjecture]{Proposition}
\newtheorem{problem}[conjecture]{Problem}
\theoremstyle{remark}
\newtheorem{remark}[conjecture]{Remark}

\newcommand{\cA}{\mathcal{A}}
\newcommand{\cC}{\mathcal{C}}
\newcommand{\barB}{\bar{B}}
\newcommand{\Ext}{\mathrm{Ext}}
\newcommand{\Hom}{\mathrm{Hom}}
\newcommand{\id}{\mathrm{id}}
\newcommand{\gr}{\operatorname{gr}}
\newcommand{\Ran}{\operatorname{Ran}}

\title{Intrinsic Homological Characterization of Chiral Koszulness:\\
A Research Memo}
\author{RED-1}
\date{March 2026}

\begin{document}
\maketitle

\section{Diagnosis of the Current State}

The monograph's definition of chiral Koszulness (Definition~\ref*{def:chiral-koszul-pair} in \texttt{chiral\_koszul\_pairs.tex}) is operationally effective but conceptually unsatisfying.  The PBW criterion (Theorem~\ref*{thm:pbw-koszulness-criterion}) states: $\cA$ is chiral Koszul if there exists a filtration $F_\bullet$ on $\cA$ such that (i) $F_\bullet$ is flat, (ii) $\gr_F \cA$ is classically Koszul, and (iii) a finiteness condition holds.  This is a \emph{sufficient} condition that reduces chiral Koszulness to classical Koszulness of the associated graded.  Three weaknesses:

\begin{enumerate}[label=(\arabic*)]
\item \textbf{Non-intrinsicality.}  The condition depends on the \emph{existence} of a good filtration.  Two chiral algebras could be quasi-isomorphic yet admit PBW filtrations whose associated gradeds have different Koszul properties.  (This does not happen for the standard examples, but the criterion does not explain \emph{why}.)

\item \textbf{Sufficient but not necessary.}  The fundamental theorem of chiral twisting morphisms (\texttt{chiral\_koszul\_pairs.tex}, Theorem, lines 272--344) gives four equivalent conditions---Koszul morphism, counit qi, unit weak equivalence, acyclicity of twisted tensor products---but all presuppose the \emph{existence} of a twisting datum $(\cA, \cC, \tau, F_\bullet)$.  This is a characterization of Koszulness \emph{relative to a given coalgebra $\cC$}, not an intrinsic property of $\cA$ alone.

\item \textbf{No obstruction theory.}  The PBW criterion gives no information about \emph{how} a non-Koszul chiral algebra fails: it says ``the spectral sequence does not degenerate at $E_2$,'' but does not identify the higher differentials as named obstructions.
\end{enumerate}

The \texttt{algebraic\_foundations.tex} file (Definition, lines 210--239) gives a cleaner formulation: $\cA$ is on the Koszul locus iff $\barB_X(\cA)$ is \emph{formal} as a factorization coalgebra, with cohomology concentrated in bar degree 1.  This is closer to intrinsic, but ``formality'' is not a checkable condition---it is the conclusion one hopes to reach.

I propose five directions for an intrinsic characterization.  Each is analyzed for precise statement, obstacles, literature support, and utility.  A ranking follows.

\section{Direction A: Ext-algebra characterization}

\textbf{Classical model.}  A connected graded algebra $A$ over a field $k$ is Koszul iff $\Ext^*_A(k,k)$ is generated in cohomological degree 1 (equivalently, $\Ext^{p,q}_A(k,k) = 0$ for $p \neq q$ in the bigrading by cohomological degree $p$ and internal degree $q$).  This is the BGS criterion \cite{BGS96}.

\textbf{Conjecture A.}
\begin{conjecture}[Chiral Ext characterization]\label{conj:ext}
Let $\cA$ be an augmented chiral algebra on a smooth projective curve $X$, with augmentation $\epsilon\colon \cA \to \omega_X$.  Define the chiral Ext algebra
\[
\Ext^*_\cA(\omega_X, \omega_X) := H^*\bigl(\mathcal{R}\!\Hom_{\cA\text{-mod}}(\omega_X, \omega_X)\bigr),
\]
computed in the derived category of chiral $\cA$-modules on $X$.  Then $\cA$ is chiral Koszul if and only if $\Ext^*_\cA(\omega_X, \omega_X)$ is generated (under the Yoneda product) in $\Ext^1$.  Equivalently, the natural map
\[
T\bigl(\Ext^1_\cA(\omega_X, \omega_X)\bigr) \longrightarrow \Ext^*_\cA(\omega_X, \omega_X)
\]
is surjective.
\end{conjecture}

\textbf{Obstacles.}
\begin{enumerate}[label=(A\arabic*)]
\item \emph{Category of chiral modules.}  The derived category $D(\cA\text{-mod})$ for a chiral algebra is not, in general, the derived category of an abelian category.  Beilinson--Drinfeld \cite{BD04} construct $\cA$-modules as $\mathcal{D}_X$-modules with compatible factorization structure, but the resulting category is not obviously abelian---it is a factorization module category, i.e., a stable $\infty$-category.  One needs to verify that $\mathcal{R}\!\Hom$ exists and computes the correct thing.  The monograph's chiral Hochschild complex (\texttt{koszul\_pair\_structure.tex}, lines 140--220) constructs this for the enveloping algebra $\cA^e$, but the self-Ext of the trivial module is a different (simpler) object.

\item \emph{Grading.}  The BGS diagonal vanishing condition requires a bigrading.  For classical graded algebras, this comes from the weight grading.  For chiral algebras, the natural ``internal'' grading is the conformal weight.  But conformal weight and bar degree do not interact as cleanly as internal degree and homological degree in the classical case: the OPE involves poles of varying order, which mix conformal weights in the bar differential.

\item \emph{Curve dependence.}  The Ext algebra depends on $X$.  For $X = \mathbb{A}^1$ (the formal disk), the chiral Ext reduces to the vertex algebra Ext and the classical criterion should apply directly.  For compact $X$ of genus $g \geq 1$, the global sections $\Gamma(X, \mathcal{R}\!\Hom)$ involve Hodge-theoretic corrections.  These are precisely the genus corrections encoded in the modular bar complex $\barB_g(\cA)$.  The Koszul locus condition should therefore be \emph{genus-dependent}: $\cA$ might be Koszul at genus 0 but fail at genus 1 (this is the admissible level phenomenon flagged at line 1664 of \texttt{bar\_cobar\_adjunction\_inversion.tex}).
\end{enumerate}

\textbf{Literature support.}  BGS \cite{BGS96}, Beilinson--Ginzburg--Soergel for the classical criterion.  Loday--Vallette \cite{LV12} for the operadic version.  Francis--Gaitsgory \cite{FG12} for factorization module categories on $\Ran(X)$.

\textbf{Utility: HIGH.}  An Ext characterization would give a concrete, checkable criterion: compute $\Ext^1$ and $\Ext^2$, check generation.  For vertex algebras at admissible levels, this would identify the precise obstruction to Koszulness.  The diagonal vanishing condition $\Ext^{p,q} = 0$ for $p \neq q$ would give a single test replacing the spectral sequence analysis.

\section{Direction B: Formality characterization}

\textbf{Classical model.}  A dg algebra $A$ is formal if it is quasi-isomorphic to its cohomology $H^*(A)$ with zero differential.  For the bar coalgebra, formality of $\barB(A)$ is equivalent to Koszulness of $A$ (this is essentially the definition in the monograph's \texttt{algebraic\_foundations.tex}, lines 210--216).

\textbf{Conjecture B.}
\begin{conjecture}[Formality characterization]\label{conj:formality}
$\cA$ is chiral Koszul if and only if $\barB_X(\cA)$, viewed as a factorization coalgebra on $\Ran(X)$, is formal: there exists a zig-zag of quasi-isomorphisms of factorization coalgebras
\[
\barB_X(\cA) \xleftarrow{\sim} \bullet \xrightarrow{\sim} H^*(\barB_X(\cA))
\]
where $H^*(\barB_X(\cA))$ is equipped with the zero differential and the coalgebra structure induced from the spectral sequence.
\end{conjecture}

\textbf{Obstacles.}
\begin{enumerate}[label=(B\arabic*)]
\item \emph{Formality of factorization coalgebras is subtle.}  In the $E_n$-algebra setting, Kontsevich formality gives formality of $E_2$-algebras but \emph{not} $E_1$-algebras (as noted in \texttt{en\_koszul\_duality.tex}, lines 96--130: ``For $E_1$-chiral algebras, the transferred $A_\infty$ structure is genuinely non-formal'').  Most chiral algebras of interest are $E_1$-chiral (ordered OPE), not $E_\infty$.  Thus Kontsevich formality does \emph{not} directly apply.

\item \emph{Curve obstructions.}  Even if $\barB(\cA)$ is formal on the formal disk, global formality on a curve $X$ of genus $g \geq 1$ faces obstructions from $H^1(X, \mathcal{O}_X) \neq 0$.  The monograph's genus-1 curvature $d_{\mathrm{fib}}^2 = \kappa \cdot \omega_1$ is precisely such an obstruction: the bar coalgebra at genus 1 is curved, hence \emph{not} formal in the naive sense.  One would need to formulate formality in the curved/filtered sense.

\item \emph{Formality vs.\ concentration.}  The monograph states (line 215--216) that formality means ``quasi-isomorphic to its cohomology, which is concentrated in bar degree 1.''  But concentration in bar degree 1 is a stronger statement than formality: it means $H^p(\barB_X(\cA)) = 0$ for $p \neq 1$.  This is the analogue of the classical statement that $\Ext^{p,q}_A(k,k) = 0$ for $p \neq q$.  Pure formality (existence of a zig-zag) does not imply concentration.  The correct conjecture should be: $\cA$ is Koszul iff $H^*(\barB_X(\cA))$ is concentrated in bar degree 1 \emph{and} the induced coalgebra structure on $H^1$ is the Koszul dual.
\end{enumerate}

\textbf{Literature support.}  Kontsevich \cite{Kon03} for $E_2$ formality.  DGMS \cite{DGMS75} for K\"ahler formality (relevant for $X$ a curve).  Positselski \cite{Positselski11} for formality in the curved setting.

\textbf{Utility: MODERATE.}  Formality is essentially a restatement of the problem.  The real content is in identifying the \emph{obstructions} to formality---which is Direction C.

\section{Direction C: $A_\infty$ obstruction theory}

\textbf{Classical model.}  By Kadeishvili's theorem, the cohomology $H^*(\barB(A))$ carries a minimal $A_\infty$-structure transferred from $\barB(A)$.  The algebra $A$ is Koszul iff this transferred structure is strict: $m_n = 0$ for all $n \geq 3$.  The higher operations $m_n$ are precisely the obstructions to formality: $m_3$ is the Massey product, $m_4$ is the quaternary Massey product, etc.

\textbf{Conjecture C.}
\begin{conjecture}[$A_\infty$ characterization]\label{conj:ainfty}
Let $\cA$ be an augmented chiral algebra on $X$.  The bar cohomology $H^*(\barB_X(\cA))$ carries a transferred $A_\infty$-factorization coalgebra structure $(m_1 = 0, m_2, m_3, m_4, \ldots)$ via the homological perturbation lemma applied to a deformation retract of $\barB_X(\cA)$ onto its cohomology.  Then:
\begin{enumerate}[label=(\roman*)]
\item $\cA$ is chiral Koszul if and only if $m_n = 0$ for all $n \geq 3$.
\item The class $[m_3] \in HH^3(\cA^!, \cA^!)$ (chiral Hochschild cohomology of the Koszul dual) is the primary obstruction to Koszulness.
\item The full sequence of obstructions $(m_3, m_4, m_5, \ldots)$ is equivalent to the sequence of higher differentials $(d_2, d_3, d_4, \ldots)$ in the PBW spectral sequence.
\end{enumerate}
\end{conjecture}

\textbf{Obstacles.}
\begin{enumerate}[label=(C\arabic*)]
\item \emph{HPL in the factorization setting.}  The homological perturbation lemma requires a deformation retract $(\barB_X(\cA), H^*(\barB_X(\cA)), p, i, h)$ where $h$ is a contracting homotopy.  For factorization coalgebras on $\Ran(X)$, the existence of $h$ is not automatic---it requires a choice of propagator on $X$.  On the formal disk, Hodge-theoretic propagators exist; on a compact curve, one needs the Arakelov Green's function or equivalent.  The propagator dependence may introduce curve-dependent obstructions not present classically.

\item \emph{$A_\infty$ vs.\ $L_\infty$.}  The monograph works with the convolution $L_\infty$-algebra $\mathrm{Conv}_\infty(\cC, \cA)$ (\texttt{algebraic\_foundations.tex}, line 686--690).  For Koszul algebras, this is formal (higher brackets vanish); for non-Koszul algebras, ``the full $L_\infty$ structure becomes essential.''  The $A_\infty$ and $L_\infty$ obstructions are related but distinct: the $A_\infty$ obstructions on $H^*(\barB(\cA))$ are Koszul-dual to the $L_\infty$ obstructions on $\cA$ itself.  One must be careful about which structure carries the obstruction.

\item \emph{Connection to the shadow tower.}  The monograph's shadow Postnikov tower $\Theta_\cA^{\leq r}$ (\texttt{higher\_genus\_modular\_koszul.tex}) is a \emph{different} kind of $A_\infty$ obstruction: it lives in the modular cyclic deformation complex and encodes genus corrections.  Conjecture C addresses genus-0 obstructions to Koszulness; the shadow tower addresses obstructions to extending the genus-0 duality to higher genus.  These are logically independent: a genus-0 Koszul algebra may still have nontrivial shadow tower (and does, whenever $\kappa \neq 0$).
\end{enumerate}

\textbf{Literature support.}  Kadeishvili \cite{Kad82} for the transferred $A_\infty$ structure.  Lu--Palmieri--Wu--Zhang \cite{LPWZ09} for the $A_\infty$ characterization of Koszulness of graded algebras.  Merkulov--Vallette for HPL in operadic settings.

\textbf{Utility: VERY HIGH.}  This gives a constructive obstruction theory: compute $m_3$, and if it is nonzero, $\cA$ is not Koszul.  The identification $m_3 \leftrightarrow d_2$ (PBW spectral sequence higher differential) would connect the abstract $A_\infty$ theory to the monograph's computational framework.  For admissible-level affine algebras, $m_3$ should be computable from the Kazhdan--Lusztig tensor product structure.

\section{Direction D: Positselski's curved perspective}

\textbf{Classical model.}  Positselski \cite{Positselski11} shows that for a CDG-algebra $(A, d, h)$ with curvature $h$ (so $d^2 = [h, -]$), the bar-cobar adjunction gives an equivalence
\[
D^{\mathrm{co}}(A\text{-mod}) \simeq D^{\mathrm{ctr}}(A^!\text{-comod})
\]
between the coderived and contraderived categories, \emph{without any Koszul hypothesis}.  This is already recorded in the monograph (\texttt{bar\_cobar\_adjunction\_inversion.tex}, Remark, lines 2187--2228).

\textbf{Conjecture D.}
\begin{conjecture}[Curved characterization]
For an augmented chiral algebra $\cA$ on $X$:
\begin{enumerate}[label=(\roman*)]
\item The bar-cobar adjunction \emph{always} gives an equivalence
\[
D^{\mathrm{co}}(\cA\text{-mod}_X) \simeq D^{\mathrm{ctr}}(\cA^!\text{-comod}_X)
\]
in the coderived/contraderived sense.
\item $\cA$ is chiral Koszul if and only if the forgetful functor $D^{\mathrm{co}}(\cA\text{-mod}_X) \to D(\cA\text{-mod}_X)$ is an equivalence, i.e., every ``absolutely acyclic'' complex is acyclic.
\item The obstruction to (ii) is measured by the curvature: the kernel of the forgetful functor is generated by the image of the curvature element $\kappa(\cA) \in HH^2_{\mathrm{cyc}}(\cA)$.
\end{enumerate}
\end{conjecture}

\textbf{Obstacles.}
\begin{enumerate}[label=(D\arabic*)]
\item \emph{Curved chiral modules.}  Positselski's theory is developed for CDG-algebras over a field.  The chiral setting requires curved factorization modules on $\Ran(X)$: the curvature $d^2 = \kappa \cdot \omega_g$ is not a scalar but a cohomology class in $H^2(\overline{\mathcal{M}}_g)$.  The passage from scalar curvature to cohomological curvature is nontrivial and not covered by existing literature.

\item \emph{Coderived vs.\ derived at genus 0.}  At genus 0, the curvature vanishes ($\kappa$ contributes only at genus $\geq 1$).  So the coderived/derived distinction is invisible at genus 0, and the curved perspective gives no new information about genus-0 Koszulness.  The interesting non-Koszul phenomena at genus 0 (admissible levels) are \emph{not} curvature phenomena---they are higher-differential phenomena in the PBW spectral sequence.  The curved perspective addresses a \emph{different} question: the interplay between genus-0 Koszulness and higher-genus curvature.

\item \emph{Positselski's hypothesis.}  The unconditional coderived equivalence requires Positselski's ``second kind'' derived functors and specific finiteness conditions (e.g., the CDG-coalgebra must be conilpotent or the CDG-algebra must be Noetherian).  Verifying these in the chiral setting is nontrivial.
\end{enumerate}

\textbf{Literature support.}  Positselski \cite{Positselski11, Positselski18}.  The monograph's Remark at lines 2187--2228 already sketches this, with the physical interpretation $\Lambda \sim \kappa(\cA)$ as cosmological constant.

\textbf{Utility: MODERATE for characterizing Koszulness, HIGH for understanding the off-Koszul regime.}  This direction does not characterize genus-0 Koszulness (which is the primary open question) but gives the correct categorical framework for the non-Koszul regime.

\section{Direction E: Derived Koszulness}

\textbf{Classical model.}  Every augmented dg-algebra $A$ admits a bar-cobar resolution $\Omega(\barB(A)) \xrightarrow{\sim} A$ in the \emph{completed} (pronilpotent) sense: the completed bar $\widehat{\barB}(A)$ always inverts under cobar.  The Koszul locus is where this completion is unnecessary.

\textbf{Conjecture E.}
\begin{conjecture}[Derived Koszulness]\label{conj:derived}
Every augmented chiral algebra $\cA$ on $X$ is ``derived Koszul'' in the following sense: the completed bar-cobar adjunction
\[
\hat{\psi}\colon \Omega_X(\widehat{\barB}_X(\cA)) \xrightarrow{\sim} \hat{\cA}
\]
is a quasi-isomorphism, where $\hat{\cA}$ and $\widehat{\barB}_X(\cA)$ denote $I$-adic completions with respect to the augmentation ideal.  The classical Koszul locus is characterized by:
\[
\cA \text{ is Koszul} \quad\Longleftrightarrow\quad \hat{\cA} \simeq \cA \text{ and } \widehat{\barB}_X(\cA) \simeq \barB_X(\cA),
\]
i.e., completion is unnecessary.
\end{conjecture}

\textbf{Obstacles.}
\begin{enumerate}[label=(E\arabic*)]
\item \emph{Completion and convergence.}  The completed bar-cobar inversion is already in the monograph (\texttt{bar\_cobar\_adjunction\_inversion.tex}, lines 405--411, citing Positselski and Gui--Li--Zeng \cite{GLZ22}).  The question is whether the completeness hypothesis is always satisfied.  For the standard examples, the $I$-adic filtration is complete and separated (Mittag-Leffler); for pathological examples, $\varprojlim^1$ terms may appear.

\item \emph{Tautology risk.}  The statement ``every algebra is derived Koszul after completion'' risks being tautological: completion is defined \emph{precisely} to make bar-cobar work.  The non-tautological content is the characterization of the Koszul locus as ``completion is unnecessary.''  But this is again a formality statement in disguise: $\widehat{\barB}(\cA) \simeq \barB(\cA)$ iff $\barB(\cA)$ is already complete, which (for conilpotent coalgebras) is automatic.  The real question is when $\hat{\cA} \simeq \cA$.

\item \emph{Non-Noetherian issues.}  Chiral algebras at admissible levels have module categories with infinite-dimensional Ext groups.  The completion $\hat{\cA}$ may not be well-behaved (e.g., the completion of a vertex algebra at an admissible level may lose the rationality that makes the representation theory tractable).
\end{enumerate}

\textbf{Literature support.}  Positselski \cite{Positselski11} for completed bar-cobar.  Gui--Li--Zeng \cite{GLZ22} (already cited in the monograph) for the filtered-complete framework.

\textbf{Utility: LOW.}  As stated, this is close to tautological.  However, the \emph{precise identification} of when completion is unnecessary could be useful if it can be formulated as a finiteness condition on $\cA$ (e.g., finite generation of the Ext algebra, or a Noetherian condition on the module category).

\section{Ranking and Recommendations}

\begin{center}
\renewcommand{\arraystretch}{1.3}
\begin{tabular}{|c|l|c|c|c|}
\hline
\textbf{Rank} & \textbf{Direction} & \textbf{Feasibility} & \textbf{Impact} & \textbf{Score} \\
\hline
1 & C: $A_\infty$ obstruction theory & 8/10 & 10/10 & \textbf{18} \\
2 & A: Ext-algebra characterization & 6/10 & 9/10 & \textbf{15} \\
3 & D: Positselski curved perspective & 7/10 & 7/10 & \textbf{14} \\
4 & B: Formality characterization & 5/10 & 6/10 & \textbf{11} \\
5 & E: Derived Koszulness & 9/10 & 3/10 & \textbf{12} \\
\hline
\end{tabular}
\end{center}

\textbf{Primary recommendation: Direction C ($A_\infty$ obstruction theory).}

The monograph already has the shadow Postnikov tower machinery for the modular ($g \geq 1$) obstruction theory.  Direction C proposes the \emph{genus-0 analogue}: the $A_\infty$ operations $(m_3, m_4, \ldots)$ on $H^*(\barB_X(\cA))$ are the obstructions to Koszulness, just as the shadow tower $(\kappa, C, Q, \ldots)$ encodes the obstructions to extending Koszul duality to higher genus.

The key theorem to prove would be:

\begin{proposition}[Target theorem]
For an augmented chiral algebra $\cA$ on a smooth curve $X$, the following are equivalent:
\begin{enumerate}[label=(\roman*)]
\item $\cA$ is chiral Koszul (in the sense of Definition~\ref*{def:chiral-koszul-pair}).
\item The minimal $A_\infty$-model of $\barB_X(\cA)$ (as a factorization coalgebra) has $m_n = 0$ for $n \geq 3$.
\item $\Ext^{p,q}_\cA(\omega_X, \omega_X) = 0$ for $p \neq q$ (diagonal vanishing).
\item The PBW spectral sequence on $\barB_X(\cA)$ degenerates at $E_2$ (existing criterion).
\end{enumerate}
Moreover, when (i)--(iv) fail, $m_3 \neq 0$ and $[m_3]$ identifies with the first nontrivial differential $d_2$ in the PBW spectral sequence.
\end{proposition}

The proof strategy for $(ii) \Leftrightarrow (iv)$ is classical (it is Kadeishvili's theorem combined with the identification of $A_\infty$ operations with Massey products, which are the transgressions of spectral sequence differentials).  The work is in establishing $(ii) \Leftrightarrow (iv)$ in the factorization setting, which requires the HPL for factorization coalgebras.

\textbf{Secondary recommendation: Direction A (Ext characterization).}  This would provide the most useful \emph{computational} tool: rather than constructing a filtration and checking classical Koszulness of the associated graded, one computes $\Ext^*_\cA(\omega_X, \omega_X)$ (which is a factorization algebra on $\Ran(X)$) and checks degree-1 generation.  The main obstacle is foundational: establishing the derived category of chiral modules with sufficient rigor.  But the monograph's chiral Hochschild cohomology apparatus (\texttt{koszul\_pair\_structure.tex}) provides the tools.

\textbf{What NOT to pursue first:}  Direction E (derived Koszulness) is essentially already in the monograph and adds no new content.  Direction B (formality) is a restatement rather than a characterization.  Direction D (curved perspective) addresses a complementary but different question (off-Koszul-locus behavior rather than Koszul characterization).

\end{document}
